%% Example data sheet
%% Feel free to modify and use this file for any purpose, under
%% either the LaTeX Project Public License or under public domain.

% Options here are passed to the article class.
% Most common options: 10pt, 11pt, 12pt

\documentclass[10pt]{../datasheet}

% Input encoding and typographical rules for English language
\usepackage[utf8]{inputenc}
\usepackage[english]{babel}
\usepackage[english]{isodate}

% tikz is used to draw images in this example, but you can
% also use \includegraphics{}.
\usepackage{tikz}
\usepackage{pgfplots}
\usepackage{circuitikz}
\usetikzlibrary{calc}
% Load Lua helper only when compiling with LuaTeX; otherwise provide a fallback
\ifdefined\directlua
  % We're in LuaTeX: load the helper
  \usepackage{luacode}
  % Netlist helper macros using Lua (requires LuaLaTeX)
\newcommand{\getnetlist}[2]{%
  \luaexec{
    local parser = dofile("../netlists/netlistparser.lua")
    local function unquote(value)
      value = tostring(value or "")
      if value:sub(1, 1) == '"' and value:sub(-1) == '"' then
        return value:sub(2, -2)
      end
      return value
    end
    local path = unquote(\luastringN{#1})
    local ref = unquote(\luastringN{#2})
    local entries = parser.parse_protel_netlist(path, ref)
    local bs = string.char(92)
    if type(entries) ~= "table" then
      tex.print(bs .. "multicolumn{2}{@{}l@{}}{Error reading netlist}" .. bs .. bs)
    elseif next(entries) == nil then
      tex.print(bs .. "multicolumn{2}{@{}l@{}}{No pins found}" .. bs .. bs)
    else
      tex.print(parser.render_pin_rows(entries))
    end
  }%
}

\else
  % Not LuaTeX: provide a harmless fallback macro so pdflatex runs cleanly
  \newcommand{\includenetlistdiagram}[2]{%
    \textbf{Netlist diagram requires compiling with LuaLaTeX (use lualatex).}%
  }
\fi

% These define global texts that are used in headers and titles.
\title{OBC v1.2 Datasheet}
\author{UTAT Space Systems}
\date{June 2026}
\revision{Revision 1}
\companylogo{\includegraphics[scale=0.022]{../figures/utat_ss.png}}
\usepackage[numbers,square]{natbib}
\usepackage{anysize}
\usepackage{graphicx} 
\marginsize{1.5cm}{1.5cm}{1cm}{1cm}
\usepackage{fancyhdr}
\renewcommand{\headrulewidth}{0pt}
\usepackage[dvipsnames]{xcolor}
\usepackage{hyperref}
\hypersetup{
    colorlinks=true,
    linkcolor=black,
    citecolor=CadetBlue,
    filecolor=CadetBlue,      
    urlcolor=CadetBlue,
}
\usepackage{multicol}
\setlength\columnsep{18pt}
\renewenvironment{abstract}
 {\par\noindent\textbf{\abstractname}\ \ignorespaces \\}
 {\par\noindent\medskip}
\renewcommand*{\thefootnote}{\fnsymbol{footnote}}
\usepackage{amsmath}


\usepackage{xcolor}
\usepackage{comment}
\usepackage{multirow}
\usepackage{float}
\usepackage{gensymb}
\usepackage{braket}
\usepackage{longtable}
\usepackage{caption}
\usepackage{longtable}
\usepackage{physics}
\usepackage{amsthm}
\usepackage{amsmath}
\usepackage{amssymb}
\usepackage{xcolor}
\usepackage{soul}
\usepackage{ marvosym }
\usepackage{wasysym}
\usepackage [autostyle, english = american]{csquotes}
\usepackage{fancyheadings}
\usepackage{enumitem}
\usepackage{multicol}


\MakeOuterQuote{"}

\newcommand{\xhat}{\hat{\vb x}}
\newcommand{\yhat}{\hat{\vb y}}
\newcommand{\zhat}{\hat{\vb z}}
\newcommand{\nhat}{\hat{\vb n}}
\newcommand{\rhat}{\hat{\vb r}}
\newcommand{\thetahat}{\hat{\vb \theta}}
\newcommand{\phihat}{\hat{\vb \phi}}


\renewcommand{\a}{\alpha}
\renewcommand{\b}{\beta}
\newcommand{\g}{\gamma}
\renewcommand{\d}{\delta}
\newcommand{\e}{\varepsilon}
\renewcommand{\l}{\lambda}
\renewcommand{\t}{\tau}
\newcommand{\w}{\omega}
\newcommand{\A}{\Alpha}
\newcommand{\B}{\Beta}
\newcommand{\G}{\Gamma}
\newcommand{\D}{\Delta}
\newcommand{\W}{\Omega}

\newcommand{\R}{\mathbb{R}}





% (netlist helper is loaded conditionally above)
\begin{document}
\maketitle

\section{High-level overview}

\begin{itemize}
  \item{Features an STM32G431RBT6 MCU}
        \begin{description}
          \item[-] ARM Cortex-M4
          \item[-] 128 kBytes of Flash
          \item[-] LQFP-64 Package
        \end{description}
  \item{HSE: 24MHz. LSE: 32.768kHz}
  \item{}
\end{itemize}

\section{Introduction}

The \textbf{datasheet} document class makes it easy to write great looking
data sheets using the LaTeX typesetting system. It follows the classic style used
by most manufacturers of electronic components.

You can download the document class from
\href{https://github.com/PetteriAimonen/latex-datasheet-template/}{latex-datasheet-template}
GitHub repository.
The repository includes this example datasheet as \textbf{example.tex} and
the document class as \textbf{datasheet.cls}.
You can build the PDF document using command \texttt{latexmk -pdf}.

% Switch to next column
\vfill\break

\begin{figure}[H]
  \centering
  \includegraphics[width=0.95\linewidth]{../figures/OBC_SS.png}
  \caption{\centering OBC v1.2}
  \label{fig:enter-label}
\end{figure}

% For wide tables, a single column layout is better. It can be switched
% page-by-page.
\onecolumn
\section{Device Overview}

Below is the block diagram that represents th high-level specification of the
On-Board Computer.

\begin{figure}[h]
  \centering
  \begin{tikzpicture}[>=latex, every node/.style={font=\small}]
    % MCU centered block
    \node[draw, rectangle, minimum width=3.2cm, minimum height=2.2cm, align=center] (mcu) {MCU\\STM32G431RBT6};

    % Tall bus block on the right (moved further right to allow label)
    \node[draw, rectangle, minimum width=1.6cm, minimum height=4cm, align=center] (bus) at ($(mcu) + (5.0,0)$) {Bus\\(I2C)};

    % I2C connection from MCU to bus
    \draw[->, thick] (mcu.east) -- node[above]{I2C (SDA / SCL)} (bus.west);

    % Single straight I2C connection
    % (SDA/SCL noted in the label on the main line)
  \end{tikzpicture}
  \caption{Block diagram: MCU and I2C bus}
\end{figure}

\section{Specifications}

TODO: cover electrical specifications such as voltage ratings and power
requirements

\section{Interconnects and Pinout}

TODO: intro

\subsection{Bus}
The table below maps the J201 connector pins found in the Protel netlist to
their net labels.

\begin{longtable}{@{}r l@{}}
  \hline
  Pin & Net \\
  \hline
  \endfirsthead
  \hline
  Pin & Net \\
  \hline
  \endhead
  \hline
  \endfoot
  \hline
  \endlastfoot
  \getnetlist{../netlists/obs_bus.net}{J201}
\end{longtable}

\subsubsection{Power}

\subsubsection{EPS connection}

\subsubsection{GPIO}

\subsubsection{GNSS}

\subsubsection{CANbus}

\subsubsection{RS-422}

\subsubsection{RS-485}

\subsubsection{PAY Debug}

\subsection{ADCS}

TODO: specify JTAG pinout and reference layout position; add the image of the
connector and the BOM part. Make sure to cover the protection circuitry in the
corresponding subsection.

\subsection{JTAG}

TODO: specify JTAG pinout and reference layout position; add the image of the
connector and the BOM part

\subsection{SRS-4}

TODO: specify the pinout and reference layout position; add the image of the
connector and BOM part

\subsection{Thermistors}

TODO: specify the pinout and reference layout position; add the image of the
connector and BOM part

TODO: describe thermistor muxes used by the OBC, use case, part numbers, and the
overall design with the opamps

\subsection{Umbillical}

TODO: specify the pinout and reference layout position.

\section{Functional Description}

\subsection{MCU}

TODO: describe some high level details about the MCU selected for the board

\subsection{External Clocks}

TODO: describe the external clocks and the reason for use, specify how they are
connected to the MCU

\subsection{External Flash Memory}

TODO: descrive the external clocks and the reason for use, specify how its
connected to the MCU.

\subsection{RTC}

TODO: describe the RTC usecase and specify how it is connected to the MCU

Real time clock (RTC) is a battery-powered external clock available to OBC for
timekeeping and synchronization purposes.

\section{Mechanical and Layout}

TODO: add a isometric view of the board with max height, dimentions and drill
holes.

TODO: add description of the polygon pours used for power dissipation

\section{Manifacturing}

\subsection{Bill of Materials}

\subsection{Test Points}

TODO: describe available test points within the design and cross reference
corresponding subsections.

\section{Bringup Log}

TODO: specify tests that have been conducted (reference to notion)

\end{document}

